\chapter{CTL over a Conversation}\label{ch:checker}

This chapter defines CTL by four recursive rules, then the states a session type implies,
which are what CTL is evaluated over, and the one notion of sameness CTL can see.
Everything in the next three chapters is stated over these definitions, and LeanFM
implements them as written.

\section{CTL}\label{sec:ctl}


A CTL formula is a statement made at one moment of a conversation about every way the
conversation can go on from there. It is built from a few atoms, which say what is true at
that moment, with $\neg$, $\wedge$, $\vee$ and $\rightarrow$, two one-step operators, and
four rules that each refer to themselves.

The atoms: $\on{\ell}$ holds on delivery of $\ell$, from the moment $\ell$ arrives until the
next message (``$\ell$ was just delivered'', not ``some time after $\ell$''; earlier versions
of LeanFM spelled it \texttt{after}, a spelling it now refuses, naming \texttt{on});
$\mathsf{stop}$ holds where the conversation may stop; $\mathsf{secret}(A, B, f)$ holds where
the step just taken carried a field ending in \texttt{\_secr} (named $f$, from $A$ to $B$,
when they are given).

\begin{definition}[CTL]\label{def:ctl}
The two steps: $\allNext\varphi$ holds at a moment when $\varphi$ holds at every moment that
can come next, and $\someNext\varphi$ when it holds at some moment that can come next.
The four recursive rules, ``else'' being ``or'':
\begin{align*}
\varphi \mathbin{\mathsf{allUntil}} \psi &\;:=\; \psi \ \text{else}\ \big(\varphi \wedge \allNext(\varphi \mathbin{\mathsf{allUntil}} \psi)\big)\\
\varphi \mathbin{\mathsf{someUntil}} \psi &\;:=\; \psi \ \text{else}\ \big(\varphi \wedge \someNext(\varphi \mathbin{\mathsf{someUntil}} \psi)\big)\\
\varphi \mathbin{\mathsf{allUnless}} \psi &\;:=\; \psi \ \text{else}\ \big(\varphi \wedge \allNext(\varphi \mathbin{\mathsf{allUnless}} \psi)\big)\\
\varphi \mathbin{\mathsf{someUnless}} \psi &\;:=\; \psi \ \text{else}\ \big(\varphi \wedge \someNext(\varphi \mathbin{\mathsf{someUnless}} \psi)\big)
\end{align*}
An \emph{until} counts only finite unfoldings of its rule: on the paths it speaks of,
$\psi$ is reached. An \emph{unless} also counts the unfolding that goes on forever:
$\varphi$ may hold all along and $\psi$ never come. A protocol satisfies a formula when the
moment before its first message does.
\end{definition}

Everything else is defined from these in one line:
\begin{align*}
\possibly\varphi &= \top \mathbin{\mathsf{someUntil}} \varphi & \inevitably\varphi &= \top \mathbin{\mathsf{allUntil}} \varphi\\
\someAlways\varphi &= \varphi \mathbin{\mathsf{someUnless}} \bot & \always\varphi &= \varphi \mathbin{\mathsf{allUnless}} \bot\\
\mathsf{never}\,\varphi &= \always\neg\varphi & \varphi\rightarrow\psi &= \neg\varphi\vee\psi
\end{align*}
$\possibly\varphi$: some path reaches $\varphi$. $\inevitably\varphi$: every path does.
$\someAlways\varphi$: some path keeps $\varphi$ forever. $\always\varphi$: every path
does, so $\varphi$ holds everywhere from here.

The shapes specifications repeat most are functions, each defined by what it unfolds to:
\begin{itemize}
\item an implication that cannot hold vacuously, written \texttt{stronglyImplies},
  \[\varphi \Rightarrow \psi = \possibly\varphi \wedge \always(\varphi\rightarrow\psi);\]
\item each outcome possible and one of them inevitable,
  \[\mathsf{outcomes}(a_1,\dots,a_n) = \possibly a_1\wedge\dots\wedge\possibly a_n\wedge
  \inevitably(a_1\vee\dots\vee a_n)\]
  (not $\inevitably a_1\vee\dots\vee\inevitably a_n$, which makes one particular outcome
  inevitable);
\item after $a$, exactly these may come next,
  \[\mathsf{follows}(a,b_1,\dots,b_n) = \always\big(a\rightarrow\someNext b_1\wedge\dots\wedge
  \someNext b_n\wedge\allNext(b_1\vee\dots\vee b_n)\big);\]
\item $\mathsf{leadsTo}(a,b) = a \Rightarrow \inevitably b$: every $a$ is answered by a $b$;
\item $\mathsf{precedes}(a,b) = \neg(\neg a\mathbin{\mathsf{someUntil}}b)$: no $b$ before $a$.
\end{itemize}
A state of the conversation has its own term, and a transition table is written with them:
\begin{align*}
(\mathsf{state}\ M) &= \on{M}\wedge\neg\mathsf{stop} &
(\mathsf{state}\ M\ \mathsf{stop}) &= \on{M}\wedge\mathsf{stop}\\
(\mathsf{transition}\ A\ B\ (C\ \mathsf{stop})) &= \mathsf{follows}(A, B, C\ \mathsf{stop}) &
(\mathsf{final}\ A) &= \mathsf{follows}(A\ \mathsf{stop}, A\ \mathsf{stop})
\end{align*}
where a bare name in a transition is a state where the conversation must go on,
$\mathsf{start}$ is the state before any message, and a final state is one where the
conversation may stop and nothing else can happen.

For readers who know the textbook terms: counting only finite unfoldings makes each until
the \emph{least} fixpoint of its rule, and counting the infinite one too makes each unless
the \emph{greatest}. Textbooks write $\varphi\mathbin{\mathsf{someUntil}}\psi$ as
$\mathsf{E}[\varphi\,\mathsf{U}\,\psi]$, the steps as $\mathsf{EX}$ and $\mathsf{AX}$, and the
derived operators as $\mathsf{EF}$, $\mathsf{AF}$, $\mathsf{EG}$, $\mathsf{AG}$
\citep{clarke1986automatic}; LeanFM reads those letters too, and writes every formula in
the words. \texttt{explain\_formula} answers any formula \emph{unfolded} into the steps
and the untils, and LeanFM's tests check its evaluator against these definitions on every
operator.

\section{The states a session type implies}

The formulas above are evaluated at moments of a conversation. A session type says what
may happen next at every point, so the moments are already in it: they are the places the
conversation can be in, and LeanFM derives them from the protocol text
(\texttt{Global.stateGraph} and \texttt{Logic.buildModel}).

A \emph{state} is what is left of the protocol to play, together with the session ids
bound so far. The first state is the whole protocol; each message the protocol allows next
leads to the state of what remains after it. Two interleavings that leave the same
remainder with the same ids bound are the same state, so a protocol with
\texttt{parallel} is no larger than its interleavings make necessary. Nested sessions are
played once in place. The states are found breadth first, at most 20{,}000 of them.

Each state remembers two things, which is all the atoms read: the message that led to it
(none, before the first message), and whether the conversation may stop there (nothing is
left that must happen). The same remainder reached by two different messages is
therefore two states. A state where the conversation may stop, or where nothing more can
happen, repeats itself forever: it \emph{stutters}. So every path goes on forever, and
$\inevitably$ and $\always$ mean what they mean in textbook CTL
\citep{clarke1986automatic}; it is also why a $\mathsf{final}$ row lists itself as what
comes next.

The later chapters use a little notation for this. The states of protocol $P$ with their
steps form its \emph{checker model} $M_P$: its states are written $n$, the first one
$n_0$, a step from $n$ to $n'$ as $n \to n'$, and the two things a state remembers as its
label $L(n) = (m, e)$, $m$ the last message and $e$ whether it may stop. The atoms read
only that label; which typestate a state belongs to has to be said with the operators of
Section~\ref{sec:ctl}, and that is the subject of Chapter~\ref{ch:faithful}.

\section{Bisimulation}

Two states are \emph{bisimilar} when they have the same label and every step one of them can
take can be matched by a step of the other to states that are again bisimilar
\citep{hennessy1985algebraic}. It is the sameness CTL can see.

\begin{definition}[Bisimulation]\label{def:bisim}
Let $M = (N, n_0, \to, L)$ and $M' = (N', n_0', \to', L')$ be total labelled graphs. A
relation $B \subseteq N \times N'$ is a \emph{bisimulation} when every $(n, n') \in B$
has $L(n) = L'(n')$, every step $n \to m$ is matched by some $n' \to' m'$ with
$(m, m') \in B$, and every step $n' \to' m'$ is matched by some $n \to m$ with
$(m, m') \in B$. $M$ and $M'$ are \emph{bisimilar}, $M \sim M'$, when some bisimulation
contains $(n_0, n_0')$. Two nodes of one model are bisimilar, $n \sim m$, when some
bisimulation of the model with itself contains $(n, m)$.
\end{definition}

\begin{theorem}[\citet{browne1988characterizing}]\label{thm:bcg}
Bisimilar nodes satisfy the same CTL formulas. Conversely, for finite models, two nodes
that satisfy the same CTL formulas are bisimilar.
\end{theorem}

The first half is a routine induction on formulas. Its consequence is the ceiling for
everything in this book: no set of CTL formulas can tell apart two systems that are
bisimilar. The best a translation into CTL can do is characterise a protocol
\emph{up to} bisimulation. The next chapter does exactly that.
